% 第9章：持续学习中四个分类器的风险表。第1阶段恒产生 (a,0)，第2阶段恒产生 (b,1)，零一损失。
% h_{uv}(a)=u, h_{uv}(b)=v。R_1(h)=I{h(a)=1}，R_2(h)=I{h(b)=0}，R_1^*=R_2^*=0。
% 平均风险 (R_1+R_2)/2，逐域超额 E_2=max(R_1,R_2)。遗忘 F_{1,2}=R_1(h^(2))-R_1(h^(1))：
% 冻结/只适应当前/兼顾新旧均以 h^(1)=h_00 为起点；“始终不学”以 h^(1)=h_10 为起点。
\begingroup
\definecolor{FigureGreenFill}{HTML}{E8EEE9}
\begin{tikzpicture}[figure picture,x=1mm,y=1mm]
\foreach \x/\w/\t in {0/24/{第2阶段结束},24/16/{$h^{(2)}$},40/18/{$(h(a),h(b))$},58/12/{$R_1$},70/12/{$R_2$},82/14/{$\overline R_2$},96/14/{$E_2$},110/16/{$F_{1,2}$}}{%
  \node[draw=FigureGrid,minimum width=\w mm,minimum height=6.4mm,inner sep=0pt,anchor=north west,font=\figurefont\footnotesize] at (\x,40) {\t};}
\foreach \y/\s/\h/\p/\ra/\rb/\av/\e/\f/\c in {
   33.6/{冻结旧模型}/{$h_{00}$}/{(0,0)}/0/1/0.5/1/0/white,
   27.2/{只适应当前}/{$h_{11}$}/{(1,1)}/1/0/0.5/1/1/FigureRedFill,
   20.8/{兼顾新旧}/{$h_{01}$}/{(0,1)}/0/0/0/0/0/FigureGreenFill,
   14.4/{始终未学好}/{$h_{10}$}/{(1,0)}/1/1/1/1/0/FigureYellowFill}{%
  \node[draw=FigureGrid,fill=\c,minimum width=24mm,minimum height=6.4mm,inner sep=0pt,anchor=north west,font=\figurefont\footnotesize] at (0,\y) {\s};
  \node[draw=FigureGrid,minimum width=16mm,minimum height=6.4mm,inner sep=0pt,anchor=north west,font=\figurefont\footnotesize] at (24,\y) {\h};
  \node[draw=FigureGrid,minimum width=18mm,minimum height=6.4mm,inner sep=0pt,anchor=north west,font=\figurefont\footnotesize] at (40,\y) {\p};
  \node[draw=FigureGrid,minimum width=12mm,minimum height=6.4mm,inner sep=0pt,anchor=north west,font=\figurefont\footnotesize] at (58,\y) {\ra};
  \node[draw=FigureGrid,minimum width=12mm,minimum height=6.4mm,inner sep=0pt,anchor=north west,font=\figurefont\footnotesize] at (70,\y) {\rb};
  \node[draw=FigureGrid,minimum width=14mm,minimum height=6.4mm,inner sep=0pt,anchor=north west,font=\figurefont\footnotesize] at (82,\y) {\av};
  \node[draw=FigureGrid,minimum width=14mm,minimum height=6.4mm,inner sep=0pt,anchor=north west,font=\figurefont\footnotesize] at (96,\y) {\e};
  \node[draw=FigureGrid,minimum width=16mm,minimum height=6.4mm,inner sep=0pt,anchor=north west,font=\figurefont\footnotesize] at (110,\y) {\f};}
\node[figure note,anchor=north west,align=left] at (0,6.5) {阶段1恒为 $(a,0)$，阶段2恒为 $(b,1)$；前三行的 $h^{(1)}=h_{00}$，末行的 $h^{(1)}=h_{10}$。};
\end{tikzpicture}
\endgroup
